\makeatletter
% mechanik.tex — die Darstellungsebene. Kennt keine Koordinaten und keine Werte.
%
% Hier steht der ganze Trick des Projekts: dieselbe Feldtabelle wird je nach
% \Modus einmal zu Formularfeldern, einmal zu festem Text, einmal zu einer
% schlichten Werteliste ohne den Originalbogen.
%
% Copyright 2026 Leif Janzik. Apache License 2.0.

% --- Werte ----------------------------------------------------------------
% \wert, \wertvon, \wertdaP, \wertoder stehen in werte.tex — dieselbe Datei,
% die die Charaktermappe einbindet. Nur so laeuft helden/<name>.tex
% unveraendert durch beide Dokumente, und zwar dauerhaft.
% werte.tex — die Werteschnittstelle. Nur diese vier Makros, sonst nichts.
%
% Sie stand wortgleich in mechanik.tex und in mappe/mechanik-mappe.tex, mit
% dem Kommentar „Gleiche Schnittstelle wie mechanik.tex, damit
% helden/<name>.tex unverändert durch beide Dokumente läuft". Genau das ist
% der Grund, warum sie nicht zweimal stehen darf: die beiden Fassungen waren
% schon auseinandergelaufen — \wertoder gab es nur in der Mappe, und \wertdaP
% war einmal einzeilig und einmal umbrochen geschrieben. Solange sie sich nur
% im Zeilenumbruch unterscheiden, fällt es nicht auf; beim nächsten Zusatz
% schon.
%
% Eingebunden mit % werte.tex — die Werteschnittstelle. Nur diese vier Makros, sonst nichts.
%
% Sie stand wortgleich in mechanik.tex und in mappe/mechanik-mappe.tex, mit
% dem Kommentar „Gleiche Schnittstelle wie mechanik.tex, damit
% helden/<name>.tex unverändert durch beide Dokumente läuft". Genau das ist
% der Grund, warum sie nicht zweimal stehen darf: die beiden Fassungen waren
% schon auseinandergelaufen — \wertoder gab es nur in der Mappe, und \wertdaP
% war einmal einzeilig und einmal umbrochen geschrieben. Solange sie sich nur
% im Zeilenumbruch unterscheiden, fällt es nicht auf; beim nächsten Zusatz
% schon.
%
% Eingebunden mit % werte.tex — die Werteschnittstelle. Nur diese vier Makros, sonst nichts.
%
% Sie stand wortgleich in mechanik.tex und in mappe/mechanik-mappe.tex, mit
% dem Kommentar „Gleiche Schnittstelle wie mechanik.tex, damit
% helden/<name>.tex unverändert durch beide Dokumente läuft". Genau das ist
% der Grund, warum sie nicht zweimal stehen darf: die beiden Fassungen waren
% schon auseinandergelaufen — \wertoder gab es nur in der Mappe, und \wertdaP
% war einmal einzeilig und einmal umbrochen geschrieben. Solange sie sich nur
% im Zeilenumbruch unterscheiden, fällt es nicht auf; beim nächsten Zusatz
% schon.
%
% Eingebunden mit \input{werte} — das Arbeitsverzeichnis ist bogen/, dafür
% sorgen die Bauskripte.
%
% \makeatletter hier selbst, nicht beim Aufrufer: \@firstoftwo braucht es,
% und eine Datei, die nur unter einer bestimmten Katcode-Lage funktioniert,
% ist eine Falle für den, der sie später woanders einbindet.
%
% Am Ende steht deshalb NICHT \makeatother, sondern die Lage des Aufrufers,
% wie sie beim Eintritt war. Der Unterschied ist einmal aufgetreten und
% kostet Suchzeit: mechanik.tex steht selbst unter \makeatletter, nach einem
% \makeatother wäre dort das @ kein Buchstabe mehr, und der Lauf endete mit
%
%   ! Undefined control sequence.
%   \feld @zeile->\endgroup \g @addto@macro\feldliste {...}
%
% — mit einem Zeilenverweis auf felder/df.tex. Die Meldung zeigt also auf
% die Feldtabelle, während der Fehler hier steht.

\expandafter\edef\csname werte@katcode\endcsname{%
  \catcode`\noexpand\@=\the\catcode`\@\relax}
\makeatletter

% \wert{feldname}{Wert} kommt aus helden/<name>.tex. Ein fehlender Wert ist
% kein Fehler, sondern ein leeres Feld — der Lauf darf daran nicht scheitern.
\newcommand{\wert}[2]{\expandafter\gdef\csname wert@#1\endcsname{#2}}

% Der Wert, oder nichts.
\newcommand{\wertvon}[1]{\ifcsname wert@#1\endcsname\csname wert@#1\endcsname\fi}

% \wertdaP{name}{dann}{sonst} — das P steht für Prädikat. Verzweigt, ohne den
% Wert selbst zu setzen.
\newcommand{\wertdaP}[1]{%
  \ifcsname wert@#1\endcsname\expandafter\@firstoftwo\else\expandafter\@secondoftwo\fi}

% Der erste gesetzte Wert aus zwei Namen, sonst leer. Damit greift die Mappe
% auf die Felder des Heldenbogens zurück: mappe_name, sonst s1_name.
\newcommand{\wertoder}[2]{\wertdaP{#1}{\wertvon{#1}}{\wertvon{#2}}}

\werte@katcode
 — das Arbeitsverzeichnis ist bogen/, dafür
% sorgen die Bauskripte.
%
% \makeatletter hier selbst, nicht beim Aufrufer: \@firstoftwo braucht es,
% und eine Datei, die nur unter einer bestimmten Katcode-Lage funktioniert,
% ist eine Falle für den, der sie später woanders einbindet.
%
% Am Ende steht deshalb NICHT \makeatother, sondern die Lage des Aufrufers,
% wie sie beim Eintritt war. Der Unterschied ist einmal aufgetreten und
% kostet Suchzeit: mechanik.tex steht selbst unter \makeatletter, nach einem
% \makeatother wäre dort das @ kein Buchstabe mehr, und der Lauf endete mit
%
%   ! Undefined control sequence.
%   \feld @zeile->\endgroup \g @addto@macro\feldliste {...}
%
% — mit einem Zeilenverweis auf felder/df.tex. Die Meldung zeigt also auf
% die Feldtabelle, während der Fehler hier steht.

\expandafter\edef\csname werte@katcode\endcsname{%
  \catcode`\noexpand\@=\the\catcode`\@\relax}
\makeatletter

% \wert{feldname}{Wert} kommt aus helden/<name>.tex. Ein fehlender Wert ist
% kein Fehler, sondern ein leeres Feld — der Lauf darf daran nicht scheitern.
\newcommand{\wert}[2]{\expandafter\gdef\csname wert@#1\endcsname{#2}}

% Der Wert, oder nichts.
\newcommand{\wertvon}[1]{\ifcsname wert@#1\endcsname\csname wert@#1\endcsname\fi}

% \wertdaP{name}{dann}{sonst} — das P steht für Prädikat. Verzweigt, ohne den
% Wert selbst zu setzen.
\newcommand{\wertdaP}[1]{%
  \ifcsname wert@#1\endcsname\expandafter\@firstoftwo\else\expandafter\@secondoftwo\fi}

% Der erste gesetzte Wert aus zwei Namen, sonst leer. Damit greift die Mappe
% auf die Felder des Heldenbogens zurück: mappe_name, sonst s1_name.
\newcommand{\wertoder}[2]{\wertdaP{#1}{\wertvon{#1}}{\wertvon{#2}}}

\werte@katcode
 — das Arbeitsverzeichnis ist bogen/, dafür
% sorgen die Bauskripte.
%
% \makeatletter hier selbst, nicht beim Aufrufer: \@firstoftwo braucht es,
% und eine Datei, die nur unter einer bestimmten Katcode-Lage funktioniert,
% ist eine Falle für den, der sie später woanders einbindet.
%
% Am Ende steht deshalb NICHT \makeatother, sondern die Lage des Aufrufers,
% wie sie beim Eintritt war. Der Unterschied ist einmal aufgetreten und
% kostet Suchzeit: mechanik.tex steht selbst unter \makeatletter, nach einem
% \makeatother wäre dort das @ kein Buchstabe mehr, und der Lauf endete mit
%
%   ! Undefined control sequence.
%   \feld @zeile->\endgroup \g @addto@macro\feldliste {...}
%
% — mit einem Zeilenverweis auf felder/df.tex. Die Meldung zeigt also auf
% die Feldtabelle, während der Fehler hier steht.

\expandafter\edef\csname werte@katcode\endcsname{%
  \catcode`\noexpand\@=\the\catcode`\@\relax}
\makeatletter

% \wert{feldname}{Wert} kommt aus helden/<name>.tex. Ein fehlender Wert ist
% kein Fehler, sondern ein leeres Feld — der Lauf darf daran nicht scheitern.
\newcommand{\wert}[2]{\expandafter\gdef\csname wert@#1\endcsname{#2}}

% Der Wert, oder nichts.
\newcommand{\wertvon}[1]{\ifcsname wert@#1\endcsname\csname wert@#1\endcsname\fi}

% \wertdaP{name}{dann}{sonst} — das P steht für Prädikat. Verzweigt, ohne den
% Wert selbst zu setzen.
\newcommand{\wertdaP}[1]{%
  \ifcsname wert@#1\endcsname\expandafter\@firstoftwo\else\expandafter\@secondoftwo\fi}

% Der erste gesetzte Wert aus zwei Namen, sonst leer. Damit greift die Mappe
% auf die Felder des Heldenbogens zurück: mappe_name, sonst s1_name.
\newcommand{\wertoder}[2]{\wertdaP{#1}{\wertvon{#1}}{\wertvon{#2}}}

\werte@katcode


% --- Feldtabelle einsammeln ------------------------------------------------
% Beim Einlesen wird nicht gesetzt, sondern gesammelt. Gesetzt wird erst je
% Seite, damit jede Seite nur ihre eigenen Felder bekommt.

% Die Argumente werden beim Einsammeln EXPANDIERT. Das ist der Grund, warum
% Tabellenzeilen mit \foreach erzeugt werden koennen: die Laufvariable (\i,
% \y) existiert nur waehrend der Schleife. Unexpandiert gespeichert waere sie
% spaeter weg, und das Feld hiesse "s4_lit_\i_probe" mit undefiniertem \i.

\newcommand{\feldliste}{}
\newcommand{\feld}[7]{%
  \begingroup
  \edef\feld@zeile{\endgroup
    \noexpand\g@addto@macro\noexpand\feldliste{%
      \noexpand\feldeintrag{#1}{#2}{#3}{#4}{#5}{#6}{#7}}}%
  \feld@zeile}

% Was mit einem gesammelten Eintrag geschieht, entscheidet der Verbraucher:
% \felderDerSeite setzt Widgets, werteliste.tex setzt Zeilen. Ohne Verbraucher
% tut ein Eintrag nichts — deshalb hier als Leerlauf vordefiniert.
\newcommand{\feldeintrag}[7]{}

% --- Feldarten -------------------------------------------------------------
% art = text | zahl | mehrzeilig
% zahl wird zentriert, mehrzeilig bekommt multiline. Alles andere wie text.
%
% Die Zusatzoptionen werden als ARGUMENT durchgereicht, nicht in einem Makro
% abgelegt. Zwei Fallgruben stecken hier drin:
%   1. \ifthenelse ist nicht expandierbar und darf nicht in der Optionsliste
%      von \TextField stehen  -> "Undefined control sequence <argument> \equal"
%   2. kvsetkeys zerlegt die Optionsliste auf Token-Ebene, BEVOR es
%      expandiert. Ein Makro, das zu "align=1" expandiert, kommt dort als ein
%      einzelnes Token an  -> "Undefined key `align=0'".
% Als Makroargument eingesetzt sind die Zeichen dagegen echte Token.

% --- Ein Feld setzen: ausfüllbare Fassung ----------------------------------
% Drei Optionen sind nicht Geschmack, sondern Pflicht:
%   bordercolor=      kein Rahmen. Das Original hat seine Linien schon.
%   backgroundcolor=  KEIN Hintergrund. Ohne das setzt hyperref /MK<</BG[1 1 1]>>
%                     und malt jedes Feld weiss ueber den Originalbogen —
%                     Linien und Beschriftungen verschwinden darunter.
%   borderwidth=0     sonst steht das Feld auf einem Rahmen.
%   bordersep=0pt     hyperrefs Vorgabe ist \Fld@bordersep = 1pt.
% charsize klein halten, sonst schneidet ein 5 mm hohes Feld unten ab.
%
% Und der eigentliche Grund fuer verschobene Rechtecke: der pdftex-Treiber
% baut Formularfelder mit \pdfstartlink, und pdfTeX blaeht JEDES Link-
% Rechteck um \pdflinkmargin auf — hyperref setzt die auf 1pt. Das sind
% 0,35 mm nach allen Seiten, auf jedem einzelnen Feld. Keine der Feldoptionen
% aendert daran etwas; es muss \pdflinkmargin selbst sein. Das Dokument hat
% keine Verweise, die Einstellung kostet also nichts.

\pdflinkmargin=0pt
%
% Mit diesen Optionen gilt: /Rect im PDF = die Werte der Feldtabelle, exakt.
% Nachpruefbar mit bau/felder-pruefen.py.

\newcommand{\feldWidget}[6]{% name x y breite hoehe zusatzoptionen
  \node[anchor=north west, inner sep=0pt, outer sep=0pt]
    at ([shift={(#2mm,-#3mm)}]current page.north west)
    {\TextField[name=#1, width=#4mm, height=#5mm,
                bordercolor=, backgroundcolor=, borderwidth=0, bordersep=0pt,
                charsize=\feldschrift, #6]{}};}

\newcommand{\setzeFeldFormular}[6]{% name x y breite hoehe art
  \ifthenelse{\equal{#6}{bild}}{}{%
  \ifthenelse{\equal{#6}{mehrzeilig}}{\feldWidget{#1}{#2}{#3}{#4}{#5}{multiline=true}}{%
  \ifthenelse{\equal{#6}{zahl}}{\feldWidget{#1}{#2}{#3}{#4}{#5}{align=1}}{%
  \feldWidget{#1}{#2}{#3}{#4}{#5}{align=0}}}}}

% --- Ein Feld setzen: vorbefuellte Fassung ---------------------------------

% Einzeilige Felder duerfen NICHT umbrechen. Eine Tabellenzelle ist hier
% 4,2 mm hoch; ein Umbruch darin laeuft in die Nachbarzeile und zerstoert
% die Tabelle. Und ein zu breiter Wert darf nicht in die Nachbarspalte
% laufen ("MU/KL/CH" in einer 12,7 mm breiten Probe-Spalte). Beides wird
% verhindert, indem der Wert erst in eine Box gesetzt und, wenn er zu breit
% ist, gestaucht wird. So bleibt jeder Wert in seiner Zelle und auf einer
% Zeile — auch wenn er dann etwas schmaler laeuft.

\newsavebox{\feldbox}

\newcommand{\einzeilig}[4]{% breite hoehe art inhalt
  \sbox\feldbox{\feldtextstil #4}%
  \ifdim\wd\feldbox>#1mm\relax
    \sbox\feldbox{\resizebox{#1mm}{!}{\usebox\feldbox}}%
  \fi
  \ifthenelse{\equal{#3}{zahl}}{%
    \parbox[t][#2mm][c]{#1mm}{\centering\usebox\feldbox}}{%
    \parbox[t][#2mm][c]{#1mm}{\raggedright\usebox\feldbox}}}

% --- Bildfelder ------------------------------------------------------------
% art=bild. Der Wert ist ein Dateipfad, nicht Text:
%   \wert{s1_bild}{helden/bilder/dorle.jpg}
% Das Bild wird in den Kasten eingepasst, Seitenverhaeltnis bleibt erhalten;
% ist es schmaler, wird es mittig gesetzt. Fehlt der Wert oder die Datei,
% bleibt der Rahmen leer — kein Abbruch, wie bei jedem anderen Feld.
%
% NUR in der vorbefuellten Fassung. AcroForm kennt kein Bildfeld; in der
% ausfuellbaren Fassung bleibt der Rahmen deshalb leer. Siehe CLAUDE.md.

\newcommand{\bildkasten}[3]{% name breite hoehe
  \ifcsname wert@#1\endcsname
    \edef\bild@pfad{\csname wert@#1\endcsname}%
    \IfFileExists{\bild@pfad}{%
      \parbox[t][#3mm][c]{#2mm}{\centering
        \includegraphics[width=#2mm, height=#3mm, keepaspectratio]{\bild@pfad}}%
    }{%
      \message{^^JHINWEIS: Bilddatei nicht gefunden: \bild@pfad^^J}%
    }%
  \fi}

\newcommand{\setzeFeldFest}[6]{%
  \node[anchor=north west, inner sep=0pt, outer sep=0pt]
    at ([shift={(#2mm,-#3mm)}]current page.north west)
    {\ifthenelse{\equal{#6}{bild}}{\bildkasten{#1}{#4}{#5}}{%
      \ifthenelse{\equal{#6}{mehrzeilig}}{%
       \parbox[t][#5mm][t]{#4mm}{\feldtextstil\raggedright\wertvon{#1}}}{%
       \einzeilig{#4}{#5}{#6}{\wertvon{#1}}}}};}

\newcommand{\feldschrift}{8pt}
\newcommand{\feldtextstil}{\fontsize{8pt}{9pt}\selectfont}

% --- Umschaltung -----------------------------------------------------------
\ifx\Modus\undefined \def\Modus{ausfuellbar}\fi

\ifthenelse{\equal{\Modus}{vorbefuellt}}{%
  \let\setzeFeld\setzeFeldFest
}{%
  \let\setzeFeld\setzeFeldFormular
}

% --- Felder einer Seite ----------------------------------------------------
% Der erste Wert der Tabellenzeile ist der LOGISCHE Bogen (1..6), nicht die
% physische Seite. Nur so heissen die Felder in beiden Quellfassungen gleich
% und dieselbe Heldendatei laeuft durch beide.

\newcommand{\aktuellerbogen}{0}

\newcommand{\feldeintragSetzen}[7]{%
  \ifnum#1=\aktuellerbogen\relax
    \setzeFeld{#2}{#3}{#4}{#5}{#6}{#7}%
  \fi}

\newcommand{\felderDerSeite}[1]{%
  \renewcommand{\aktuellerbogen}{#1}%
  \begin{tikzpicture}[remember picture, overlay]
    \let\feldeintrag\feldeintragSetzen
    \feldliste
    \ifmessen\messgitterinhalt\fi
  \end{tikzpicture}}

% --- Das Messgitter --------------------------------------------------------
% Der Kern der Machbarkeit: ohne Gitter muessten ~200 Positionen geraten
% werden, mit ihm werden sie abgelesen. Bleibt dauerhaft im Projekt.
% Alle Koordinaten in Millimeter von der LINKEN OBEREN Papierecke.

\newif\ifmessen
\ifx\Messen\undefined\messenfalse\else\messentrue\fi

\newcommand{\messgitterinhalt}{%
  \foreach \x in {0,10,...,210}{%
    \draw[cyan!40, line width=0.1pt]
      ([shift={(\x mm,0)}]current page.north west) --
      ([shift={(\x mm,0)}]current page.south west);
    \node[cyan, font=\tiny, anchor=north]
      at ([shift={(\x mm,-2mm)}]current page.north west) {\x};}
  \foreach \y in {0,10,...,290}{%
    \draw[cyan!40, line width=0.1pt]
      ([shift={(0,-\y mm)}]current page.north west) --
      ([shift={(0,-\y mm)}]current page.north east);
    \node[cyan, font=\tiny, anchor=west]
      at ([shift={(2mm,-\y mm)}]current page.north west) {\y};}
  % Feinlinien alle 5 mm, unbeschriftet — fuer das genaue Ablesen
  \foreach \x in {5,15,...,205}{%
    \draw[cyan!15, line width=0.05pt]
      ([shift={(\x mm,0)}]current page.north west) --
      ([shift={(\x mm,0)}]current page.south west);}
  \foreach \y in {5,15,...,285}{%
    \draw[cyan!15, line width=0.05pt]
      ([shift={(0,-\y mm)}]current page.north west) --
      ([shift={(0,-\y mm)}]current page.north east);}}

% --- Leere Boegen weglassen -----------------------------------------------
% Fuer einen unmagischen Helden ist Bogen 5 (Zauber) eine leere Seite, fuer
% einen Weltlichen ebenso Bogen 4 (Liturgien). Drei Wege, sie loszuwerden;
% der erste, der etwas sagt, gewinnt:
%
%   1. \bogennur{1,2,3,6}       in der Heldendatei - nur diese Boegen
%   2. \bogenweglassen{4,5}     in der Heldendatei - alle ausser diesen
%   3. Bauoption -OhneLeere     automatisch: ein Bogen ohne einen einzigen
%                               gesetzten Wert faellt weg
%
% Weg 3 wirkt NUR in der vorbefuellten Fassung. In der ausfuellbaren gibt es
% keine Werte; dort wuerde er alle Boegen wegwerfen. Er wird deshalb dort
% ignoriert, mit Hinweis im Log.
%
% Die Eigenschaftsleiste MU..KK zaehlt bei Weg 3 NICHT als Inhalt. Sonst
% haette jeder Bogen Werte, sobald die acht Eigenschaften gesetzt sind, und
% die Automatik liefe leer - genau der Fall, der bei Dorle aufgetreten waere.

\newcommand{\zaehltnicht}[1]{\expandafter\gdef\csname zaehltnicht@#1\endcsname{}}

\newcommand{\werteZaehlen}{%
  \begingroup
    \def\feldeintrag##1##2##3##4##5##6##7{%
      \ifcsname wert@##2\endcsname
        \ifcsname zaehltnicht@##2\endcsname\else
          \expandafter\gdef\csname bogenhat@##1\endcsname{}%
        \fi
      \fi}%
    \feldliste
  \endgroup}

\newcommand{\bogenauswahl}{}      % leer = keine explizite Auswahl
\newcommand{\bogennur}[1]{\gdef\bogenauswahl{,#1,}\gdef\bogenmodus{nur}}
\newcommand{\bogenweglassen}[1]{\gdef\bogenauswahl{,#1,}\gdef\bogenmodus{ohne}}
\newcommand{\bogenmodus}{}

\newif\ifbogenzeigen

\newcommand{\bogenpruefen}[1]{%
  \bogenzeigentrue
  \ifx\bogenmodus\@empty
    % keine explizite Auswahl: nur die Automatik kann noch etwas wegnehmen
    \ifx\OhneLeere\undefined\else
      \ifthenelse{\equal{\Modus}{ausfuellbar}}{%
        \message{^^JHINWEIS: -OhneLeere wirkt nur in der vorbefuellten Fassung.^^J}%
      }{%
        \ifcsname bogenhat@#1\endcsname\else\bogenzeigenfalse\fi
      }%
    \fi
  \else
    % \in@ durchsucht sein zweites Argument auf TOKEN-Ebene. Uebergibt man
    % dort das Makro \bogenauswahl, sieht \in@ ein einzelnes Token und findet
    % nie etwas — die Auswahl bliebe wirkungslos, ohne Fehlermeldung.
    % Deshalb erst expandieren, dann suchen.
    \begingroup\edef\bogen@x{\endgroup\noexpand\in@{,#1,}{\bogenauswahl}}\bogen@x
    \ifthenelse{\equal{\bogenmodus}{nur}}{%
      \ifin@\else\bogenzeigenfalse\fi
    }{%
      \ifin@\bogenzeigenfalse\fi
    }%
  \fi}

% Faellt JEDER Bogen weg, entsteht ein Dokument ohne Seiten — pdflatex
% schreibt dann gar kein PDF, und das sieht nach einem Absturz aus statt nach
% einer Auswahl. Fuer diesen Fall bleibt Bogen 1 stehen.
\newif\ifbogennotfall

\newcommand{\bogenAuswahlPruefen}{%
  \bogennotfalltrue
  \bogenpruefen{1}\ifbogenzeigen\bogennotfallfalse\fi
  \bogenpruefen{2}\ifbogenzeigen\bogennotfallfalse\fi
  \bogenpruefen{3}\ifbogenzeigen\bogennotfallfalse\fi
  \bogenpruefen{4}\ifbogenzeigen\bogennotfallfalse\fi
  \bogenpruefen{5}\ifbogenzeigen\bogennotfallfalse\fi
  \bogenpruefen{6}\ifbogenzeigen\bogennotfallfalse\fi
  \ifbogennotfall
    \message{^^JHINWEIS: kein Bogen hat Inhalt. Bogen 1 wird trotzdem
      ausgegeben, damit ein PDF entsteht.^^J}%
  \fi}

% --- Eine Bogenseite ausgeben ---------------------------------------------
% Originalseite als Grund, Felder absolut darueber. Nicht nachbauen.

\newcommand{\bogenseite}[1]{%
  \bogenpruefen{#1}%
  \ifbogennotfall
    \ifnum#1=1 \bogenzeigentrue\else\bogenzeigenfalse\fi
  \fi
  \ifbogenzeigen
    \includepdf[pages={\quellseite{#1}},
                picturecommand={\felderDerSeite{#1}}]{\quelldatei}%
  \else
    \message{^^JBogen #1 weggelassen (leer oder abgewaehlt).^^J}%
  \fi}
\makeatother
